\otazka{Kyberválka (doplňkové téma)}{Definice, historie, současné trendy. Identifikace aktérů, problémy s přisouzením kyberútoků, možnost odstrašení v kyberprostoru. Koncept netwars a jeho použití v kyberkonfliktu.}{BSSb1152}

\subsection*{Kyberválka -- stručněji}
U kyberválky je nejdůležitější nepoužít slovo válka na každý hack. Odpověď musí rozlišit kriminalitu, špionáž, sabotáž, informační operace, kyberkonflikt a skutečnou kyberválku.

\textbf{1. Definice a hranice pojmu}
\begin{itemize}
  \item Kyberválka je úzký pojem pro použití kyberprostředků v kontextu ozbrojeného konfliktu nebo k účinkům srovnatelným s vojenským útokem. Ne každý státem podporovaný kyberútok je válka.
  \item Kyberoperace může být špionáž, příprava bojiště, sabotáž, nátlak, propaganda nebo kriminální zisk. Proto se vždy ptám na aktéra, cíl, dopad, kontext a práh použití síly.
\end{itemize}

\textbf{2. Historie a příklady}
\begin{itemize}
  \item Estonsko 2007 ukazuje masivní DDoS proti státu závislému na digitálních službách. Dopad byl hlavně dostupnost webů, bank a veřejných služeb, ne fyzické zničení.
  \item Gruzie 2008 ukazuje propojení kyberútoků s konvenční vojenskou operací. DDoS a defacementy doprovázely ruskou vojenskou akci a měly narušit komunikaci a informační obraz.
  \item Viasat 2022 ukazuje útok na satelitní komunikaci na začátku ruské invaze na Ukrajinu. Malware narušil modemy KA-SAT, dopadl na ukrajinskou komunikaci i některé evropské uživatele.
\end{itemize}

\textbf{3. Aktéři a atribuce}
\begin{itemize}
  \item Aktéři mohou být státy, zpravodajské služby, armády, proxy skupiny, kyberkriminální skupiny, hacktivisté, insideři a soukromí dodavatelé.
  \item Atribuce je těžká, protože útočník používá proxy infrastrukturu, kompromitované servery, falešné stopy a sdílené nástroje. Přisouzení proto není jen technické, ale i zpravodajské a politické.
\end{itemize}

\textbf{4. Odstrašení}
\begin{itemize}
  \item Odstrašení trestem znamená hrozbu nákladů: sankce, trestní stíhání, diplomatická reakce, veřejná atribuce nebo protiopatření.
  \item Odstrašení odepřením znamená, že útok se nevyplatí, protože obrana, segmentace, zálohy a obnova sníží dopad.
  \item V kyberprostoru je odstrašení těžší kvůli anonymitě, nízkým nákladům útoku, nejistému prahu reakce a rozdílu mezi státními a nestátními aktéry.
\end{itemize}

\textbf{5. Netwars}
\begin{itemize}
  \item Netwar označuje síťově organizovaný konflikt, kde se propojuje informační působení, mobilizace komunit, propaganda, DDoS, hacktivismus a tlak na veřejné mínění.
  \item Příklad NoName057(16): proruská hacktivistická síť vyhlašuje cíle, mobilizuje podporovatele a používá DDoS proti webům států podporujících Ukrajinu. Dopad je často symbolický a reputační, ale politicky viditelný.
\end{itemize}

\subsection*{Kostra odpovědi}
\begin{enumerate}
  \item Přesně odlišit incident, kriminalitu, špionáž, sabotáž, informační operaci, kyberkonflikt, kyberobranu a kyberválku.
  \item Definovat kyberválku úzce: ne každý státem podporovaný útok je kyberválka.
  \item Vysvětlit historický vývoj přes konkrétní příklady a jejich význam.
  \item Popsat současné trendy vždy s případem: APT31, Volt Typhoon, Salt Typhoon, APT28/TP-Link, NoName057(16), AI, dodavatelé a soukromý sektor.
  \item Vysvětlit aktéry, atribuci, odstrašení a základní mezinárodněprávní rámec.
  \item Ukázat netwar na proruských hacktivistických kampaních.
\end{enumerate}

\subsection*{Přesné rozlišení pojmů}
U této otázky je nejdůležitější nezačít tím, že každý velký kyberútok označíš za kyberválku. \textbf{Kybernetický incident} je událost, která narušuje nebo může narušit bezpečnost systému, služby nebo dat. Může být technickou chybou, omylem, kriminalitou nebo státní operací. \textbf{Kyberkriminalita} je trestná činnost, například ransomware, phishing, podvod, neoprávněný přístup nebo vydírání. \textbf{Kyberšpionáž} je získávání informací, přístupů a dokumentů; může být státní i nestátní a typicky zůstává pod prahem ozbrojeného útoku. \textbf{Kybernetická sabotáž} míří na narušení, poškození nebo vyřazení systému či procesu. \textbf{Informační operace} míří na vnímání, rozhodování a důvěru veřejnosti nebo institucí.

\textbf{Kybernetický konflikt} je širší pojem pro soutěžení a střet v kyberprostoru: špionáž, sabotáž, nátlak, DDoS, vlivové operace, úniky dat, příprava přístupu nebo operace během války. \textbf{Kyberobrana} je vojenská a státně-obranná dimenze působení v kyberprostoru. \textbf{Kyberválka} je užší pojem: použití kybernetických operací státem nebo jiným politickým aktérem jako součást ozbrojeného konfliktu nebo k dosažení účinků srovnatelných se závažným vojenským útokem.

Rozhodující je aktér, účel, rozsah, účinky a vztah k širšímu konfliktu. APT31 proti MZV ČR, Salt Typhoon nebo Volt Typhoon jsou velmi významné státní nebo státem podporované operace, ale bez dalšího je nelze automaticky označit za ozbrojený útok nebo kyberválku. U zkoušky stačí říct: "Kyberválka je úzký pojem, většina škodlivých kyberoperací zůstává v šedé zóně pod prahem války."

\subsection*{Kyberprostor, vrstvy a typy účinků}
Kyberprostor není jen "internet". Pro zkoušku je užitečné ho rozdělit na několik vrstev. \textbf{Fyzická vrstva} jsou kabely, servery, datová centra, satelity, koncová zařízení a průmyslová zařízení. \textbf{Logická/protokolová vrstva} jsou adresace, směrování, protokoly, domény, autentizace a síťová pravidla. \textbf{Datová nebo sémantická vrstva} je obsah: databáze, e-maily, dokumenty, konfigurace, účty. \textbf{Kognitivní vrstva} je lidské vnímání, rozhodování, důvěra a informační působení.

Kyberoperace mohou mít různé účinky. \textbf{Logický účinek} je změna dat, účtů, konfigurace nebo dostupnosti služby. \textbf{Fyzický účinek} nastává, když zásah do systému ovlivní reálný proces, například turbínu, tlak, signalizaci, nemocniční provoz nebo průmyslovou linku. \textbf{Kognitivní účinek} míří na důvěru, rozhodování, paniku nebo politickou legitimitu. Proto je důležitá triáda CIA - důvěrnost, integrita, dostupnost - ale u konfliktu je potřeba přidat i odolnost, obnovu a dopad na společnost.

\priklad{Když útočník ukradne diplomatické e-maily, jde hlavně o důvěrnost a politický účinek. Když změní konfiguraci průmyslového řídicího systému, jde o integritu a možný fyzický účinek. Když shodí online služby státu před volbami a současně šíří narativ o neschopnosti vlády, kombinuje dostupnost a kognitivní účinek.}

Typy kyberoperací se dají vysvětlit podle toho, co chtějí způsobit. \textbf{Akvizice informací} znamená špionáž a krádež dat. \textbf{Zveřejnění nebo šíření informací} znamená úniky, propagandu a dezinformace. \textbf{Disrupce} znamená narušení provozu, například DDoS nebo vyřazení části služby. \textbf{Destrukce} znamená mazání dat, wipery nebo výjimečně fyzické poškození zařízení. Tím se dá pěkně ukázat, že kyberválka není jen "hackování", ale práce s účinky na rozhodování, provoz a schopnost státu jednat.

\pozor{Kyberválku nepleť s elektronickým bojem. Elektronický boj pracuje s elektromagnetickým spektrem, rušením, radiolokací a signály. Může vojensky souviset s kyberoperacemi, ale není to totéž.}

\subsection*{Historický vývoj přes příklady}
Historii neříkej jako seznam let, ale jako několik případů, na kterých ukážeš vývoj: nejdřív DDoS a politický tlak, potom zásah do kritické infrastruktury, destruktivní malware s globálními dopady a nakonec kyberoperace přímo navázaná na konvenční válku.

\priklad{\textbf{Případová studie - Estonsko 2007.} Kdy a kontext: na jaře 2007 se Estonsko dostalo do sporu s Ruskem kvůli přesunu sovětského památníku v Tallinnu. Současně začaly vlny DDoS útoků proti estonským vládním webům, parlamentu, ministerstvům, médiím, bankám, poskytovatelům internetu a DNS službám. Co se technicky dělo: útočníci zahlcovali servery provozem, používali botnety i dobrovolnickou mobilizaci přes fóra a návody, takže služba nebyla dostupná běžným uživatelům. Cíl: nešlo o fyzické zničení infrastruktury, ale o ochromení digitálního státu, politický tlak, reputační dopad a vytvoření dojmu, že stát neumí chránit vlastní online služby. Dopad: výpadky veřejných webů, bankovnictví a médií; zasažena byla důvěra v digitalizovaný stát. Reakce: Estonsko zapojilo domácí správce sítí, mezinárodní technickou pomoc a incident se stal impulzem pro posílení kyberobrany NATO a rozvoj Tallinnského centra CCDCOE. Zkouškový význam: dobrý příklad kyberkonfliktu a netwar, ale opatrně s pojmem kyberválka, protože přímá právní přičitatelnost ruskému státu nebyla jednoduše prokázána.}

\priklad{\textbf{Případová studie - ukrajinská energetika 2015/2016.} Kdy a kdo: 23. prosince 2015 zasáhl útok ukrajinské distribuční energetické společnosti; bezpečnostní komunita jej běžně spojuje se skupinou Sandworm a ruskými zájmy, ukrajinská SBU vinila Rusko. Co bylo cílem: narušit distribuci elektřiny a ukázat schopnost zasáhnout základní službu během širšího rusko-ukrajinského konfliktu. Jak to proběhlo: útočníci se do prostředí dostali dlouhodobě předem, mimo jiné přes spear phishing a malware BlackEnergy, získali přístupy do provozních systémů, vzdáleně ovládali rozvodny, vypínali jističe, mazali systémy pomocí KillDisk a zároveň zahltili call centra, aby zákazníci nemohli hlásit výpadky. Dopad: výpadky elektřiny pro velký počet odběratelů a nutnost ruční obnovy provozu. V roce 2016 následoval sofistikovanější útok spojený s malwarem Industroyer/CrashOverride, který byl navržen pro komunikaci s průmyslovými protokoly. Reakce: manuální provoz, obnova konfigurací, oddělování IT a OT, posilování detekce a mezinárodní sdílení poznatků. Zkouškový význam: tady už jasně ukážeš přechod z IT do OT a fyzický dopad na základní službu.}

\priklad{\textbf{Případová studie - NotPetya 2017.} Kdy a kdo: v červnu 2017 se z Ukrajiny rozšířil destruktivní malware NotPetya; Spojené království a USA jej veřejně připsaly ruskému státu/ruské armádě. Cíl: podle veřejných atribucí destabilizovat Ukrajinu, zejména její instituce a ekonomické prostředí. Jak to technicky začalo: malware byl rozšířen přes kompromitovaný aktualizační mechanismus ukrajinského účetního softwaru M.E.Doc, který používala řada firem v Ukrajině. Jak se šířil: uvnitř sítí využíval techniky pro laterální pohyb a credential abuse, takže rychle zasáhl celé podnikové prostředí. Důležité: tvářil se jako ransomware, ale prakticky šlo o wiper; zaplacení výkupného nemělo reálně vést k obnově dat. Dopad: ochromení ukrajinských organizací a rozsáhlé globální škody u firem jako Maersk, Merck nebo FedEx/TNT; problém se přelil přes dodavatelské a podnikové vazby mimo původní cíl. Reakce: veřejná atribuce západních států, diplomatické odsouzení, posílení debaty o sankcích, kyberpojištění a odpovědnosti za státní destruktivní operace. Zkouškový význam: destruktivní operace maskovaná jako kriminalita a příklad nekontrolovaných vedlejších dopadů.}

\priklad{\textbf{Případová studie - Viasat KA-SAT 2022.} Kdy a kdo: útok proběhl 24. února 2022, tedy v den zahájení ruské invaze proti Ukrajině; EU a spojenci jej veřejně spojili s Ruskou federací. Cíl: narušit satelitní komunikační službu KA-SAT, kterou využívali uživatelé na Ukrajině i v dalších evropských zemích. Jak to proběhlo podle Viasatu: útočník zneužil chybnou konfiguraci nebo kompromitovaný přístup ve správní části sítě, získal možnost posílat příkazy zákaznickým modemům a poškodil jejich funkčnost. Dopad: část zákazníků ztratila satelitní připojení, bylo nutné měnit nebo obnovovat koncová zařízení a dopady se projevily i mimo Ukrajinu, například u civilních evropských uživatelů. Reakce: technická obnova služby a výměna zařízení, veřejné odsouzení a politická atribuce ze strany EU a partnerů. Zkouškový význam: kyberoperace časově sladěná s konvenční vojenskou operací a ukázka, že komerční infrastruktura může mít civilně-vojenský význam a přeshraniční dopady.}

\priklad{\textbf{Případová studie - APT31 proti MZV ČR.} Kdy a kdo: česká vláda v květnu 2025 veřejně atribuovala dlouhodobou škodlivou kampaň proti jedné z neutajovaných sítí Ministerstva zahraničních věcí ČR skupině APT31 spojované s čínskou zpravodajskou službou; podle NÚKIB aktivita probíhala minimálně od roku 2022. Cíl: získávat diplomatické informace a dlouhodobý přístup k infrastruktuře ministerstva, ne způsobit viditelný výpadek. Jak to chápat: jde o kyberšpionáž a působení pod prahem války; technické detaily nemusejí být veřejné, protože stát chrání zdroje a metody. Dopad: bezpečnostní a diplomatický problém, nikoli fyzická škoda. Reakce: veřejná politická atribuce, zapojení NÚKIB, vlády, MZV, zpravodajských informací a spojenců. Zkouškový význam: ideální případ pro vysvětlení rozdílu mezi technickou atribucí, zpravodajským hodnocením, politickou atribucí a právní přičitatelností.}

\subsection*{Současné trendy}
\textbf{Působení pod prahem války.} Dnešní kyberkonflikt se často odehrává v šedé zóně: špionáž, krádeže přístupů, DDoS, vlivové operace, zveřejňování odcizených dat, kompromitace účtů a příprava přístupu pro budoucí konflikt. Cílem nemusí být okamžitě něco vypnout; často stačí získat informace, oslabit důvěru nebo držet možnost budoucího nátlaku.

\aktualni{\textbf{APT31 proti MZV ČR.} Česká vláda v květnu 2025 veřejně atribuovala škodlivou kampaň skupině APT31 spojované s čínskou zpravodajskou službou. Podle NÚKIB cílila minimálně od roku 2022 na jednu z neutajovaných sítí Ministerstva zahraničních věcí ČR. Příklad vysvětluje kyberšpionáž, veřejnou politickou atribuci a šedou zónu: jde o závažnou státní operaci proti infrastruktuře státu, ale ne automaticky o kyberválku.}

\textbf{Pre-positioning v kritické infrastruktuře.} Útočník nemusí infrastrukturu hned vypnout. Může si tiše připravovat přístup, mapovat závislosti, získávat účty a držet perzistenci pro budoucí politickou nebo vojenskou krizi.

\aktualni{\textbf{Volt Typhoon.} CISA a partneři popsali čínským státem podporovaného aktéra, který si udržoval dlouhodobý přístup v prostředí americké kritické infrastruktury. Typické byly living off the land techniky, legitimní účty, silná operační bezpečnost a dlouhodobá perzistence. Význam případu není v prokázané destrukci, ale v přípravě možnosti budoucího narušení energetiky, vody, dopravy, komunikací nebo dalších služeb v krizi.}

\textbf{Telekomunikace jako strategický cíl.} Telekomunikační sítě mají zpravodajskou hodnotu, protože přes jednoho poskytovatele lze sledovat provozní údaje, správní přístupy, komunikaci vybraných osob a vztahy mezi institucemi. Routery a síťová zařízení jsou proto strategické cíle, ne jen technická periferie.

\aktualni{\textbf{Salt Typhoon.} NÚKIB se v srpnu 2025 připojil ke společnému upozornění amerických a dalších partnerů na čínské státem podporované aktéry, jejichž aktivity částečně odpovídají označení Salt Typhoon. Cílem byly zejména telekomunikační sítě, dopravní infrastruktura, ubytovací zařízení a vojenské systémy. Případ ukazuje strategickou hodnotu páteřních routerů, provider edge/customer edge zařízení, důvěryhodných propojení a dlouhodobého přístupu pro špionáž.}

\textbf{Zneužití infrastruktury třetích stran.} Zdrojová IP adresa neprokazuje identitu útočníka. Státní i kriminální aktéři mohou používat domácí routery, malé firemní routery, VPS, cloudové účty nebo botnety jako prostředníky. Tím kryjí původ operace a komplikují atribuci.

\aktualni{\textbf{APT28 a kompromitované routery TP-Link.} NÚKIB v dubnu 2026 informoval, že ruský státní aktér APT28 podle zpráv USA a FBI vybudoval globální síť kompromitovaných SOHO routerů. U zranitelného modelu TP-Link TL-WR841N měl zneužít známou zranitelnost, měnit DHCP/DNS nastavení a přes podvržené DNS odpovědi získávat hesla, tokeny a další citlivé informace. Příklad ukazuje, proč samotná IP adresa nestačí, proč jsou okrajová zařízení strategická a proč jsou výchozí hesla a EOL zařízení bezpečnostní problém.}

\textbf{Hacktivismus a proxy aktéři.} DDoS se používá jako nástroj politického tlaku, mediální viditelnosti a swarmingového zahlcení. U proruských skupin je potřeba rozlišovat autentický hacktivismus, státní toleranci, státní koordinaci a právní přičitatelnost. To, že skupina tvrdí fyzický dopad, ještě neznamená, že se skutečně stal.

\aktualni{\textbf{NoName057(16) a další proruské skupiny.} NÚKIB v prosinci 2025 upozornil spolu s partnery na proruské hacktivistické útoky proti kritické infrastruktuře. U NoName057(16) NÚKIB evidoval od roku 2023 desítky incidentů proti českým institucím a firmám, všechny jako DDoS. Upozornění zároveň řešilo i oportunistické útoky na špatně zabezpečená OT/ICS zařízení přes vzdálený přístup. Příklad se hodí pro hacktivismus, proxy aktéry, netwar a rozdíl mezi reálným dopadem a propagandistickým tvrzením.}

\textbf{AI v kyberkonfliktu.} AI realisticky zvyšuje rychlost a kvalitu phishingu, překladů, jazykové lokalizace, deepfake audia a videa, analýzy odcizených dat, průzkumu a skriptování. Stejně tak pomáhá obraně: korelace logů, detekce anomálií, triage incidentů a tvorba pravidel. Není vhodné tvrdit, že už existuje běžná plně autonomní kyberválka bez lidské kontroly.

\textbf{Dodavatelské řetězce a poskytovatelé.} Kompromitace jednoho důvěryhodného dodavatele, cloudu, telekomunikačního operátora, správního nástroje nebo aktualizačního mechanismu může otevřít cestu k mnoha obětem. Proto se útoky přesouvají od jednotlivých počítačů k identitám, tokenům, správním účtům, API, cloudovým oprávněním a třetím stranám.

\textbf{Soukromý sektor jako aktér konfliktu.} Válka na Ukrajině ukazuje význam cloudových společností, satelitních poskytovatelů, threat-intelligence firem, výrobců technologií, dobrovolnických skupin a veřejně-soukromých partnerství. Soukromé firmy nejsou jen pasivní oběti; mohou poskytovat obrannou infrastrukturu, analýzu hrozeb, obnovu služeb i politicky citlivé služby dvojího užití.

\subsection*{Aktéři}
\begin{itemize}
  \item \textbf{Státy a jejich vojenské či zpravodajské složky} mají strategickou motivaci, dlouhodobé zdroje a cíle jako špionáž, příprava bojiště, sabotáž nebo podpora vojenských operací.
  \item \textbf{Státem podporované APT skupiny} mohou být formálně oddělené od státu, ale jejich cíle často odpovídají státním zájmům.
  \item \textbf{Proxy aktéři} jsou skupiny, které stát toleruje, podporuje nebo využívá, aby si udržel popiratelnost.
  \item \textbf{Kyberkriminální skupiny} sledují hlavně zisk, například ransomware nebo krádeže dat, ale jejich infrastruktura se může potkat se státními zájmy.
  \item \textbf{Hacktivisté} jsou motivovaní ideologií, solidaritou nebo politickým protestem a typicky cílí na viditelnost.
  \item \textbf{Teroristické organizace} běžně používají internet k propagandě, náboru, komunikaci a financování; mnohem vzácnější je kybernetický útok, který by sám naplnil znaky teroristického násilí.
  \item \textbf{Soukromé firmy} provozují cloud, satelitní komunikaci, bezpečnostní služby, threat intelligence a někdy i nástroje dvojího užití.
  \item \textbf{Jednotlivci a dobrovolnické skupiny} mají menší zdroje, ale v asymetrickém prostředí mohou způsobit viditelný dopad nebo se zapojit do širších kampaní.
\end{itemize}
U aktérů se vždy ptej na motivaci, zdroje, typické cíle a vztah ke státu. Stejný technický útok může mít jiný význam, když ho provede kriminální skupina kvůli výkupnému, hacktivisté kvůli publicitě nebo státní aktér kvůli přípravě na budoucí konflikt.

\subsection*{Atribuce kyberútoků}
Atribuce je proces, kterým se kybernetická operace přisuzuje konkrétnímu aktérovi. Není to jen otázka IP adresy. IP adresa ukazuje, odkud přišlo spojení, nikoli nutně kdo útok provedl.

\begin{longtable}{p{0.17\linewidth}p{0.43\linewidth}p{0.32\linewidth}}
\toprule
\textbf{Úroveň atribuce} & \textbf{Co se zkoumá} & \textbf{Příklad}\\
\midrule
Technická & Malware, infrastruktura, TTP, logy, účty, indikátory kompromitace. & APT28 a kompromitované routery TP-Link.\\
Operační & Výběr cílů, pracovní rytmus, opakující se postupy skupiny. & Dlouhodobé kampaně APT aktérů.\\
Strategická & Motiv, geopolitický kontext a komu operace prospívá. & APT31 proti MZV ČR.\\
Zpravodajská & Neveřejné informace, zdroje a metody, které stát nemusí zveřejnit. & Veřejně se zveřejní jen část podkladů.\\
Politická & Rozhodnutí státu útok veřejně přisoudit. & Česká veřejná atribuce APT31 v roce 2025.\\
Právní & Přičitatelnost podle mezinárodního práva. & Vyšší práh než technická pravděpodobnost.\\
\bottomrule
\end{longtable}

Hlavní problémy jsou VPN a proxy, botnety, kompromitovaná zařízení, infrastruktura třetích osob, sdílený nebo odcizený malware, falešné stopy, změněné časové údaje, false flag operace, rozdíl mezi pravděpodobností a jistotou a ochrana zpravodajských zdrojů. Technické přiřazení není totéž jako právní přičitatelnost. Veřejná atribuce je politické rozhodnutí a stát často nezveřejní všechny důkazy, protože by prozradil zpravodajské zdroje a metody.

\aktualni{\textbf{APT31 a politická atribuce.} Česká atribuce APT31 je dobrý příklad konečné politické roviny. Student má říct: nejprve technická a zpravodajská práce, potom vyhodnocení strategického kontextu a nakonec politické rozhodnutí vlády, co veřejně říct a jaký diplomatický signál vyslat.}

\subsection*{Odstrašení v kyberprostoru}
Odstrašení znamená přesvědčit protivníka, že útok se mu nevyplatí, protože nebude úspěšný nebo za něj ponese nepřijatelnou cenu. V kyberprostoru je to složitější než u klasického vojenského odstrašení: atribuce je nejistá, útoky často zůstávají pod prahem války, aktéři jsou různí a náklady na útok mohou být nízké.

Rozlišují se tři základní logiky:
\begin{itemize}
  \item \textbf{Odstrašení trestem} - sankce, trestní stíhání, diplomatické kroky, zveřejnění atribuce nebo jiná odveta.
  \item \textbf{Odstrašení odepřením} - útok je technicky obtížný a nepřinese očekávaný výsledek díky segmentaci, MFA, zálohám, monitoringu a řízení přístupů.
  \item \textbf{Odolnost} - i úspěšný útok nezpůsobí dlouhodobý strategický efekt, protože stát nebo organizace službu obnoví.
\end{itemize}

K tomu přidej \textbf{kolektivní reakci}: společné atribuce, koordinované sankce, EU, NATO, sdílení informací a diplomatický tlak. Vedle toho existuje \textbf{rozbíjení infrastruktury útočníků}: policejní, zpravodajské nebo technické operace proti botnetům, malware infrastruktuře, serverům, doménám nebo kryptoměnovým tokům.

\aktualni{\textbf{Jak spojit příklady s odstrašením.} Volt Typhoon se hodí hlavně k odstrašení odepřením a odolnosti: segmentace, hardening, detekce living off the land, ochrana identit a schopnost fungovat i při kompromitaci. APT31 se hodí k odstrašení trestem a kolektivní reakci: veřejná atribuce, diplomacie a politický tlak. DDoS kampaně NoName057(16) se hodí k otázce, zda rychlá obrana a obnova odeberou útočníkovi propagandistický efekt.}

\pozor{Těžký kyberútok může vytvářet riziko eskalace, ale není přesné tvrdit, že téměř vždy přeroste do fyzického násilí. U zkoušky mluv o riziku eskalace, ne o automatickém následku.}

Na příkladu nemocnice nebo energetiky je rozdíl dobře vidět. Odstrašení trestem říká: když zaútočíš, vystavuješ se sankcím, trestnímu stíhání nebo diplomatické reakci. Odstrašení odepřením říká: útok se ti nevyplatí, protože máme segmentaci, monitoring a řízené přístupy. Odolnost říká: i když něco zasáhneš, nezískáš strategický efekt, protože službu obnovíme.

\subsection*{Mezinárodní právo}
Mezinárodní právo platí i v kyberprostoru. Sporná není samotná použitelnost práva, ale jeho aplikace na konkrétní technické situace. \textbf{Suverenita} chrání státní autoritu nad územím, infrastrukturou a veřejnými funkcemi. \textbf{Zákaz vměšování} chrání rozhodovací autonomii státu, například volby nebo výkon veřejné moci. \textbf{Zákaz použití síly} může být relevantní, pokud kyberoperace svými účinky odpovídá použití síly, například rozsáhlým fyzickým škodám nebo paralyzování kritické infrastruktury.

\textbf{Ozbrojený útok} je ještě vyšší práh a může založit právo na sebeobranu. \textbf{Retorze} je nepřátelský, ale dovolený krok, například diplomatické vyhoštění nebo sankce. \textbf{Protiopatření} by jinak byla protiprávní, ale za splnění podmínek mohou být dočasnou reakcí na předchozí mezinárodně protiprávní jednání jiného státu. \textbf{Odpovědnost státu} vyžaduje přičitatelnost jednání státu; politická atribuce a právní přičitatelnost nejsou totéž.

V ozbrojeném konfliktu se použije i mezinárodní humanitární právo. Kyberoperace musí respektovat rozlišování civilních a vojenských cílů, proporcionalitu, vojenskou nezbytnost a ochranu civilistů. Tallinn Manual je odborná nezávazná studie, ne smlouva a ne právně závazný předpis. NATO uvádí, že závažný kybernetický útok může vést k úvahám o článku 5, ale vždy politicky a případ od případu; neplatí automaticky "kyberútok = článek 5".

\subsection*{Netwars}
Netwar je konflikt vedený síťově organizovanými aktéry, kteří využívají decentralizovanou komunikaci, informace a koordinované působení proti hierarchicky organizovanému protivníkovi. Koncept vytvořili John Arquilla a David Ronfeldt. Netwar není jen hacking; může zahrnovat propagandu, mobilizaci, dezinformace, psychologické působení, hacktivismus, úniky dat, koordinované kampaně a tlak na veřejné mínění.

Rozdíl je tento: \textbf{kyberválka} se týká hlavně využití kybernetických operací pro státní nebo vojenské cíle. \textbf{Netwar} se týká způsobu organizace a vedení konfliktu. Nemusí probíhat pouze v kyberprostoru, ale kyberprostor mu dává rychlou koordinaci, propagandu, nábor, informační operace a swarming.

V kyberkonfliktu se netwars projevují tak, že vedle technických útoků běží informační operace. Útočník nejen shodí službu, ale zároveň šíří narativ, že stát je neschopný. Nejen ukradne e-maily, ale zveřejní je ve správný politický moment. Nejen napadne infrastrukturu, ale posiluje chaos na sociálních sítích.

Síťové konflikty mohou mít různé organizační podoby. \textbf{Řetězová síť} předává komunikaci nebo zboží po linii. \textbf{Hvězdicová síť} má centrální uzel, kolem kterého se koordinují ostatní. \textbf{Vícekanálová neboli all-channel síť} umožňuje komunikaci mezi všemi uzly a je nejméně závislá na jednom centru. \textbf{Swarming} je koordinované rojení většího počtu menších aktérů nebo útoků proti jednomu či více cílům z různých směrů.

Netwar je dobré vysvětlit jako konflikt sítí proti hierarchii. Tradiční stát rozhoduje pomaleji, přes instituce, procedury a odpovědnosti. Síťově organizovaní aktéři mohou jednat rychle, decentralizovaně a často bez jasného velitele. To je výhoda u kampaní, kde se kombinuje únik dat, DDoS, memy, dezinformace, dobrovolnické skupiny, kriminální infrastruktura a mediální tlak. Obrana pak není jen technická, ale i komunikační a politická: stát musí chránit systémy, vysvětlovat situaci veřejnosti a nenechat si vnutit útočníkův narativ.

\aktualni{\textbf{Proruské hacktivistické kampaně jako netwar.} Skupiny typu NoName057(16) zveřejňují seznamy cílů, mobilizují dobrovolníky přes sociální sítě a komunikační platformy, provádějí DDoS a okamžitě publikují propagandistické výsledky. Technicky jde často "jen" o dostupnost, ale politicky o tlak na veřejné mínění a reputaci státu. To je swarming: více menších útoků a mediálních signálů působí současně na technické systémy i psychologii publika.}

\subsection*{Banka příkladů ke zkoušce}
Tuhle část ber jako rychlé karty. U každého případu si zkus říct: \emph{kdy, kdo, cíl, technika, dopad, reakce, proč je to důležité}.
\begin{itemize}
  \item \textbf{Estonsko 2007 - politický DDoS a digitální stát.} Po sporu o sovětský památník byly zahlcovány weby vlády, parlamentu, médií, bank a poskytovatelů. Technicky šlo hlavně o DDoS, politicky o nátlak a reputační zásah. Reakce vedla k posílení mezinárodní kyberobranné spolupráce a často se zmiňuje v souvislosti s NATO CCDCOE. Neříkej zkratkovitě "Rusko hacknulo Estonsko"; správně: šlo o politicky motivované útoky s problematickou atribucí.
  \item \textbf{Ukrajinská energetika 2015/2016 - IT útok s fyzickým dopadem.} Útočníci se dostali do sítí distribučních společností, získali přístupy, ovládali provozní systémy, vypínali rozvodny, mazali stanice a komplikovali komunikaci se zákazníky. Dopadem byly výpadky elektřiny a ruční obnova. Reakce: segmentace IT/OT, obnova manuálních postupů, sdílení poznatků o ICS hrozbách. Význam: ukázka, že kyberoperace může narušit základní službu, i když nejde o tanky nebo bomby.
  \item \textbf{NotPetya 2017 - wiper maskovaný jako ransomware.} Šíření začalo přes kompromitovanou aktualizaci ukrajinského účetního softwaru M.E.Doc. Malware se uvnitř sítí rychle šířil a ničil systémy; výkupné bylo spíš krycí legenda než reálná cesta k obnově. Dopad zasáhl Ukrajinu i globální firmy, například logistiku, farmaceutický sektor a dopravu. Reakce: veřejná atribuce Rusku ze strany USA, UK a dalších, debata o sankcích, kyberpojištění a odpovědnosti za státní destruktivní operace.
  \item \textbf{Viasat KA-SAT 2022 - kyberoperace v okamžiku invaze.} V den ruské invaze proti Ukrajině byla narušena satelitní komunikační služba KA-SAT. Podle Viasatu útočník zneužil správní přístup a poslal destruktivní příkazy koncovým modemům, což vedlo k výpadkům a nutnosti obnovy nebo výměny zařízení. Dopad byl přeshraniční a dotkl se civilních uživatelů mimo Ukrajinu. Reakce: technická obnova a veřejná atribuce/odsouzení ze strany EU a spojenců.
  \item \textbf{APT31 proti MZV ČR - špionáž, ne automaticky kyberválka.} Česká vláda v roce 2025 veřejně atribuovala dlouhodobou kampaň proti neutajované síti MZV skupině APT31 spojované s čínskou zpravodajskou službou. Cílem bylo získávání informací a přístupů, ne fyzická destrukce. Reakce: veřejná atribuce, diplomatický signál a spolupráce se spojenci. Význam: rozdíl mezi technickou stopou, zpravodajským hodnocením, politickou atribucí a právní přičitatelností.
  \item \textbf{Volt Typhoon - připravený přístup bez okamžité destrukce.} CISA a partneři popsali čínského státem podporovaného aktéra v prostředí americké KI. Nešlo o prokázané vypnutí elektrárny nebo vody, ale o dlouhodobé udržování přístupu, living off the land a zneužívání legitimních účtů. Cíl: možnost budoucího narušení během krize. Význam: pre-positioning, odstrašení odepřením a odolnost.
  \item \textbf{NoName057(16) - hacktivismus a netwar.} Skupina zveřejňuje cíle, mobilizuje dobrovolníky, provádí DDoS a okamžitě komunikuje propagandistické výsledky. Dopad bývá často hlavně dostupnostní a reputační, ne fyzická destrukce. Reakce: filtrace DDoS, rychlá obnova, komunikace s veřejností a spolupráce s NÚKIB/partnery. Význam: swarming, síťový konflikt a rozdíl mezi reálným dopadem a propagandou.
\end{itemize}

\zkouska{Nejdřív odliš incident, kriminalitu, špionáž, sabotáž, informační operaci, kyberkonflikt, kyberobranu a kyberválku. Kyberválka je úzký pojem: kyberoperace státu nebo politického aktéra jako součást ozbrojeného konfliktu nebo s účinky srovnatelnými se závažným vojenským útokem. Ne každý státní útok je kyberválka. U příkladů neříkej jen název: vždy řekni kdy, aktéra nebo míru atribuce, cíl, techniku, dopad a reakci. Estonsko 2007 = DDoS a politický tlak na digitální stát s problematickou atribucí. Ukrajinská energetika 2015/2016 = průnik přes IT do OT, vypínání rozvoden, mazání systémů a ruční obnova. NotPetya 2017 = wiper přes ukrajinský software M.E.Doc, globální vedlejší dopady a veřejná atribuce Rusku. Viasat 2022 = satelitní komunikace narušená v den invaze, dopad i mimo Ukrajinu a reakce EU/spojenců. APT31 proti MZV = špionáž a politická atribuce, ne automaticky kyberválka. Atribuce má technickou, operační, strategickou, zpravodajskou, politickou a právní rovinu. Odstrašení funguje přes trest, odepření, odolnost, kolektivní reakci a rozbíjení infrastruktury útočníků. Netwar je síťový konflikt: decentralizace, propaganda, DDoS, mobilizace dobrovolníků a swarming.}

\zdroje{Přednáškový materiál \emph{3.) Kyberválka (LONG)}; přiložené PDF a BSS poznámky; Tallinn Manual jako akademický rámec; NATO a EU dokumenty k odpovědnému chování a kolektivní reakci; CCDCOE k Estonsku; CISA k ukrajinské energetice a Volt Typhoon; Viasat a Rada EU k útoku KA-SAT; UK/USA k NotPetya; NÚKIB k APT31, APT28/TP-Link, Salt Typhoon a NoName057(16); NATO CCDCOE k Ukrajině.}

\newpage

